\section{Setting}
\label{sec:problem_setting}

In this work, we consider RL problems with both multiple agents and partial observability. All the agents share the goal of maximising the same discounted sum of rewards $R_t$.  While no agent can observe the underlying Markov state $s_t$, each agent receives a private observation $o^a_t$ correlated to $s_t$.  In each time-step, the agents select an \emph{environment action} $u \in U$ that affects the environment, and a \emph{communication action} $m \in M$ that is observed by other agents but has no direct impact on the environment or reward.
We are interested in such settings because it is only when multiple agents and partial observability coexist that agents have the incentive to communicate. As no communication protocol is given a priori, the agents must develop and agree upon such a protocol to solve the task.

Since protocols are mappings from action-observation histories to sequences of messages, the space of protocols is extremely high-dimensional. Automatically discovering effective protocols in this space remains an elusive challenge.  In particular, the difficulty of exploring this space of protocols is exacerbated by the need for agents to coordinate the sending and interpreting of messages.  For example, if one agent sends a useful message to another agent, it will only receive a positive reward if the receiving agent correctly interprets and acts upon that message.  If it does not, the sender will be discouraged from sending that message again.  Hence,  positive rewards are sparse, arising only when sending and interpreting are properly coordinated, which is hard to discover via random exploration.

We focus on settings with \emph{centralised learning} but \emph{decentralised execution}.  In other words, communication between agents is not restricted during learning, which is performed by a centralised algorithm; however, during execution of the learned policies, the agents can communicate only via the limited-bandwidth channel.  While not all real-world problems can be solved in this way, a great many can, e.g., when training a group of robots on a simulator.  Centralised planning and decentralised execution is also a standard paradigm for multi-agent planning \cite{Oliehoek08JAIR,kraemer2016multi}. 



